\documentclass[11pt]{beamer}
% =====================
% Thème
% =====================
\usetheme{CambridgeUS}
\usecolortheme{default}
% =====================
% Langue et encodage
% =====================
\usepackage[french]{babel}
\usepackage[T1]{fontenc}
\usepackage{lmodern}
% =====================
% Maths et figures
% =====================
\usepackage{amsmath, amssymb}
\usepackage{graphicx}
\usepackage[most]{tcolorbox}
% =====================
% Infos du document
% =====================
\title{Modélisation de la sévérité et de la fréquence des accidents
routiers}
\author[Équipe 15]{%
BKHACH Ayoub\\
BEGGAR Rayane\\
SABRI Yassir\\
MARSOUK Ayoub
}

\institute[ENSIMAG]{%
ENSIMAG\\[0.5em]
\textbf{Encadrant :} Christophe Dutang
}

\date{\today}

% =====================
% Début du document
% =====================
\begin{document}

% ---------- Slide de titre ----------
\begin{frame}
  \titlepage
\end{frame}

% ---------- Plan ----------
\begin{frame}{Plan}
  \tableofcontents
\end{frame}

% =====================
% Section 1
% =====================
\section{Introduction}

\begin{frame}{Contexte}
\begin{itemize}
  \item Présentation du problème
  \item Motivation scientifique
  \item Objectifs du travail
\end{itemize}
\end{frame}

% =====================
% Section 2
% =====================
\section{Score de sévérité}

\begin{frame}{Définition du Score de sévérité et Modèle linéaire :}
\begin{columns}

\column{0.48\textwidth}
\vspace{-0.5cm}
\[
\tcboxmath{
S = \frac{9T + 3H + B}{13} \times 100
}
\]
\column{0.48\textwidth}
\vspace{-0.1cm}
\[
\tcboxmath{
S = \beta_0 + \sum_{j=1}^{p} \beta_j X_j
}
\]

\end{columns}
\textbf{Exemples des variables explicatives : } Infrastructure (agglomération, type de route), Environnement (luminosité, conditions météorologiques), Accident (type de collision, nombre de véhicules impliqués et leurs types)

\vspace{+0.2cm}
\small
\centering
\begin{tabular}{l c c}
\hline
\textbf{Condition atmosphérique} & \textbf{Coefficient} & \textbf{p-valeur} \\
\hline
Brouillard / fumée        & $-4.73$ & $0.0025$ \\
Couvert                  & $-3.95$ & $0.0016$ \\
Éblouissant               & $-2.36$ & $0.0774$ \\
Neige / grêle             & $+0.67$ & $0.7429$ \\
Conditions normales       & $-4.60$ & $< 10^{-4}$ \\
Pluie forte               & $-5.04$ & $< 10^{-4}$ \\
Pluie légère              & $-5.56$ & $< 10^{-5}$ \\
Vent fort / tempête       & $+0.58$ & $0.7569$ \\
\hline
\end{tabular}
\end{frame}



\begin{frame}{Modèle linéaire  simplifié :}
\begin{itemize}
    \item \textbf{Idée :} Uilisation de la fonction (step) et dimunier le nombre de variable selon le sens logique des mots et la p-valeur
\end{itemize}

\vspace{0.3cm}

\textbf{Luminosité :}
\[
\begin{array}{ccl}
\text{Plein jour} & \longrightarrow & \text{Avec lumière} \\
\text{Nuit avec éclairage allumé} & \longrightarrow & \text{Avec lumière} \\
\text{Nuit avec éclairage non allumé} & \longrightarrow & \text{Sans lumière} \\
\text{Nuit sans éclairage} & \longrightarrow & \text{Sans lumière}
\end{array}
\]

\vspace{0.4cm}

\textbf{Conditions atmosphériques :}
\[
\begin{array}{ccl}
\text{Normale, Couvert} & \longrightarrow & \text{Normal} \\
\text{Pluie légère, Éblouissant} & \longrightarrow & \text{Défavorable} \\
\text{Brouillard, Neige/Grêle, Vent fort, Pluie forte} & \longrightarrow & \text{Sévère}
\end{array}
\]

\end{frame}

\section{Les regroupements à l’aide de méthode de partitionnement}
\begin{frame}{Répartition des clusters (avec modèle complet) : }
\vspace{-0.1cm}
\begin{figure}
    \centering
    \includegraphics[width=1\linewidth]{répartition non simplifié .png}
\end{figure}
\end{frame}

\begin{frame}{Répartition des clusters (avec modèle simplifié et sans sévérité)}
\vspace{-0.25cm}
\begin{figure}
        \centering
        \includegraphics[width=1\linewidth]{répartition simplifié.png}
\end{figure}
\end{frame}
\begin{frame}{Répartition des clusters (avec modèle simplifié et avec sévérité)}
\vspace{-0.25cm}
\begin{figure}
    \centering
    \includegraphics[width=1\linewidth]{répartition simplifié avec sévérité.png}
    \caption{Enter Caption}
    \label{fig:placeholder}
\end{figure}
\end{frame}
% =====================
% freéquences
% =====================
\section{fréquences et modèle quasi poisson}
\begin{frame}{fréquences par mois : }
\begin{figure}
    \centering
    \includegraphics[width=1\linewidth]{fréquences par mois.png}
\end{figure}
\end{frame}

\begin{frame}{fréquences par departement : }
\begin{figure}
    \centering
    \includegraphics[width=1\linewidth]{fréquences par département.png}
    \caption{Enter Caption}
    \label{fig:placeholder}
\end{figure}
\end{frame}
% =====================
% loi quasi pois
% =====================
\begin{frame}{Modèle en loi de poisson : }
\begin{itemize}
  \item \textbf{Intercept (554)}
  \begin{itemize}
    \item Niveau moyen mensuel d’accidents à Paris (situation de référence)
  \end{itemize}

  \item \textbf{Proportion d’accidents de nuit}
  \begin{itemize}
    \item Effet négatif marqué
    \item $\Rightarrow$ Diminution du nombre d’accidents mensuels
  \end{itemize}

  \item \textbf{Météo défavorable} \; ($\exp(\beta)\approx 4{,}9$)
  \begin{itemize}
    \item Effet fortement positif
    \item $\Rightarrow$ Forte augmentation du risque d’accident
  \end{itemize}
  \item \textbf{Collisions fortes} \; ($\exp(\beta)\approx 4{,}8$)
  \begin{itemize}
    \item Effet positif majeur
    \item $\Rightarrow$ Accidents plus fréquents
  \end{itemize}
\end{itemize}
\end{frame}
% =====================
% Conclusion
% =====================
\section{Conclusion}
\begin{frame}{Conclusion}
\begin{itemize}
  \item Bilan
  \item Limites
  \item Perspectives
\end{itemize}
\end{frame}
% ---------- Slide final ----------
\begin{frame}
\centering
\Large Merci pour votre attention
\end{frame}
\end{document}